% GENERATED FILE. Edit canonical lessons, then run npm run generate:books.
% canonical-source-hash: fnv1a64:6b33e4092508552a
% canonical-lessons: AR-C320-ilam, AR-C320-wakala, AR-C320-masdar, AR-C320-afada, AR-C320-awdaha

\chapter{The media}
\label{ch:ar-the-media}
\hlchaptermodality{320}

\begin{chapteropening}
\textbf{By the end of this chapter:} \emph{I can say the media, an agency, a source, to report and to clarify.}

Complete the last lesson of chapter 320: I can say the media, an agency, a source, to report and to clarify.
\end{chapteropening}

\section[\texorpdfstring{\ar{إعلام} (iʿlām)}{إعلام (iʿlām)}]{\ar{إعلام} (iʿlām) — the media; informing}

\label{lesson:AR-C320-ilam}



\begin{quote}
Before the new one: say the Arabic for to announce a death, then the Arabic for mourning.
\end{quote}

\subsection*{You'll want to know: \ar{إعلام}}
\textbf{\ar{إعلام}} — \emph{iʿlām} — \textquotedblleft{}the media; informing\textquotedblright{}.

\begin{quote}
\ar{تكلم} \ar{الإعلام} \ar{كثيرا} \ar{عن} \ar{الحادث}.
\end{quote}

\begin{quote}
\emph{takallama al-iʿlām kathīran ʿan al-ḥādith.}
\end{quote}

\begin{quote}
The media talked a lot about the accident.
\end{quote}

\begin{cousinweb}
\textbf{\ar{إعلام}} (\emph{iʿlām}) is built on the native Arabic root \emph{ʿ-l-m} (Lane), \textquotedblleft{}to know\textquotedblright{}, the root of \emph{ʿilm}, knowledge. \emph{Iʿlām} is the verbal noun of Form IV \emph{aʿlama}, \textquotedblleft{}to let someone know\textquotedblright{}.
\end{cousinweb}

\begin{grammarlens}[title={What you've built}]
One more word for telling what happened.
\end{grammarlens}

\subsection*{Your turn}
\begin{itemize}
\raggedright
  \item \emph{Say it:} \emph{iʿlām}
  \item \emph{Say it:} \emph{iʿlām}, once more
  \item \emph{Say it:} \emph{takallama al-iʿlām kathīran ʿan al-ḥādith.} — The media talked a lot about the accident.
  \item \emph{Recall:} say the Arabic for to announce a death, then the Arabic for mourning, then say \emph{iʿlām} again
\end{itemize}

\begin{tcolorbox}[breakable,colback=teal!4,colframe=teal!35!black,title={Before you move on}]
What does \ar{إعلام} mean? (\textquotedblleft{}The media; informing\textquotedblright{}.) Say it once more.
\end{tcolorbox}
\section[\texorpdfstring{\ar{وكالة} (wakāla)}{وكالة (wakāla)}]{\ar{وكالة} (wakāla) — an agency}

\label{lesson:AR-C320-wakala}



\begin{quote}
Before the new one: say the Arabic for mourning, then the Arabic for the media.
\end{quote}

\subsection*{You'll want to know: \ar{وكالة}}
\textbf{\ar{وكالة}} — \emph{wakāla} — \textquotedblleft{}an agency\textquotedblright{}.

\begin{quote}
\ar{نشرت} \ar{وكالة} \ar{الأنباء} \ar{الخبر} \ar{صباحا}.
\end{quote}

\begin{quote}
\emph{nasharat wakālat al-anbāʾ al-khabar ṣabāḥan.}
\end{quote}

\begin{quote}
The news agency put out the story in the morning.
\end{quote}

\ar{وكالة} \ar{سفر}: a travel agency.

\begin{cousinweb}
\textbf{\ar{وكالة}} (\emph{wakāla}) is built on the native Arabic root \emph{w-k-l} (Lane), \textquotedblleft{}to entrust an affair to someone\textquotedblright{}. A \emph{wakīl} acts for others; a \emph{wakāla} is the office that does.
\end{cousinweb}

\begin{grammarlens}[title={What you've built}]
One more word for telling what happened.
\end{grammarlens}

\subsection*{Your turn}
\begin{itemize}
\raggedright
  \item \emph{Say it:} \emph{wakāla}
  \item \emph{Say it:} \emph{wakāla}, once more
  \item \emph{Say it:} \emph{nasharat wakālat al-anbāʾ al-khabar ṣabāḥan.} — The news agency put out the story in the morning.
  \item \emph{Recall:} say the Arabic for mourning, then the Arabic for the media, then say \emph{wakāla} again
\end{itemize}

\begin{tcolorbox}[breakable,colback=teal!4,colframe=teal!35!black,title={Before you move on}]
What does \ar{وكالة} mean? (\textquotedblleft{}An agency\textquotedblright{}.) Say it once more.
\end{tcolorbox}
\section[\texorpdfstring{\ar{مصدر} (maṣdar)}{مصدر (maṣdar)}]{\ar{مصدر} (maṣdar) — a source}

\label{lesson:AR-C320-masdar}



\begin{quote}
Before the new one: say the Arabic for the media, then the Arabic for an agency.
\end{quote}

\subsection*{You'll want to know: \ar{مصدر}}
\textbf{\ar{مصدر}} — \emph{maṣdar} — \textquotedblleft{}a source\textquotedblright{}.

\begin{quote}
\ar{قال} \ar{مصدر} \ar{في} \ar{الشرطة} \ar{إن} \ar{اللص} \ar{هرب}.
\end{quote}

\begin{quote}
\emph{qāla maṣdar fī ash-shurṭa inna al-liṣṣ haraba.}
\end{quote}

\begin{quote}
A police source said the thief had escaped.
\end{quote}

Plural: \ar{مصادر}, \emph{maṣādir}.

In grammar the \ar{مصدر} is the verbal noun.

\begin{cousinweb}
\textbf{\ar{مصدر}} (\emph{maṣdar}) is built on the native Arabic root \emph{ṣ-d-r} (Lane), \textquotedblleft{}to go back from the watering place\textquotedblright{}. \emph{Maṣdar} is the place a thing comes out of: a source.
\end{cousinweb}

\begin{grammarlens}[title={What you've built}]
One more word for telling what happened.
\end{grammarlens}

\subsection*{Your turn}
\begin{itemize}
\raggedright
  \item \emph{Say it:} \emph{maṣdar}
  \item \emph{Say it:} \emph{maṣdar}, once more
  \item \emph{Say it:} \emph{qāla maṣdar fī ash-shurṭa inna al-liṣṣ haraba.} — A police source said the thief had escaped.
  \item \emph{Recall:} say the Arabic for the media, then the Arabic for an agency, then say \emph{maṣdar} again
\end{itemize}

\begin{tcolorbox}[breakable,colback=teal!4,colframe=teal!35!black,title={Before you move on}]
What does \ar{مصدر} mean? (\textquotedblleft{}A source\textquotedblright{}.) Say it once more.
\end{tcolorbox}
\section[\texorpdfstring{\ar{أفاد} (afāda)}{أفاد (afāda)}]{\ar{أفاد} (afāda) — to give news; to be of use}

\label{lesson:AR-C320-afada}



\begin{quote}
Before the new one: say the Arabic for an agency, then the Arabic for a source.
\end{quote}

\subsection*{You'll want to know: \ar{أفاد}}
\textbf{\ar{أفاد}} — \emph{afāda} — \textquotedblleft{}to give news; to be of use\textquotedblright{}.

\begin{quote}
\ar{أفادت} \ar{الشرطة} \ar{بأن} \ar{الطريق} \ar{مغلق}.
\end{quote}

\begin{quote}
\emph{afādat ash-shurṭa bi-anna aṭ-ṭarīq mughlaq.}
\end{quote}

\begin{quote}
The police reported that the road was closed.
\end{quote}

Present: \ar{يفيد}, \emph{yufīd}.

In the news: \ar{أفاد} \ar{مراسلنا} \ar{بأن} …, our correspondent reports that …. Elsewhere: \ar{هذا} \ar{يفيد}, this helps.

\begin{cousinweb}
\textbf{\ar{أفاد}} (\emph{afāda}) is built on the native Arabic root \emph{f-y-d} (Lane), where \emph{fāʾida} is a gain or benefit. Form IV \emph{afāda} is \textquotedblleft{}to give someone a benefit\textquotedblright{}, and so to give him information.
\end{cousinweb}

\begin{grammarlens}[title={What you've built}]
One more word for telling what happened.
\end{grammarlens}

\subsection*{Your turn}
\begin{itemize}
\raggedright
  \item \emph{Say it:} \emph{afāda}
  \item \emph{Say it:} \emph{afāda}, once more
  \item \emph{Say it:} \emph{afādat ash-shurṭa bi-anna aṭ-ṭarīq mughlaq.} — The police reported that the road was closed.
  \item \emph{Recall:} say the Arabic for an agency, then the Arabic for a source, then say \emph{afāda} again
\end{itemize}

\begin{tcolorbox}[breakable,colback=teal!4,colframe=teal!35!black,title={Before you move on}]
What does \ar{أفاد} mean? (\textquotedblleft{}To give news; to be of use\textquotedblright{}.) Say it once more.
\end{tcolorbox}
\section[\texorpdfstring{\ar{أوضح} (awḍaḥa)}{أوضح (awḍaḥa)}]{\ar{أوضح} (awḍaḥa) — to clarify, to make clear}

\label{lesson:AR-C320-awdaha}



\begin{quote}
Before the new one: say the Arabic for a source, then the Arabic for to give news.
\end{quote}

\subsection*{You'll want to know: \ar{أوضح}}
\textbf{\ar{أوضح}} — \emph{awḍaḥa} — \textquotedblleft{}to clarify, to make clear\textquotedblright{}.

\begin{quote}
\ar{أوضح} \ar{المعلم} \ar{الدرس} \ar{بمثال}.
\end{quote}

\begin{quote}
\emph{awḍaḥa al-muʿallim ad-dars bi-mithāl.}
\end{quote}

\begin{quote}
The teacher made the lesson clear with an example.
\end{quote}

Present: \ar{يوضح}, \emph{yūḍiḥ}.

\begin{cousinweb}
\textbf{\ar{أوضح}} (\emph{awḍaḥa}) is built on the native Arabic root \emph{w-ḍ-ḥ} (Lane), \textquotedblleft{}to be plain, clear\textquotedblright{}, the root of \emph{wāḍiḥ}, clear. It is Form IV, \textquotedblleft{}to make clear\textquotedblright{}.
\end{cousinweb}

\begin{grammarlens}[title={What you've built}]
One more word for telling what happened.
\end{grammarlens}

\subsection*{Your turn}
\begin{itemize}
\raggedright
  \item \emph{Say it:} \emph{awḍaḥa}
  \item \emph{Say it:} \emph{awḍaḥa}, once more
  \item \emph{Say it:} \emph{awḍaḥa al-muʿallim ad-dars bi-mithāl.} — The teacher made the lesson clear with an example.
  \item \emph{Recall:} say the Arabic for a source, then the Arabic for to give news, then say \emph{awḍaḥa} again
\end{itemize}

\begin{tcolorbox}[breakable,colback=teal!4,colframe=teal!35!black,title={Before you move on}]
What does \ar{أوضح} mean? (\textquotedblleft{}To clarify, to make clear\textquotedblright{}.) Say it once more.
\end{tcolorbox}
